\section{Sparsification on the instance}\label{S:sparsification}

At $h=\tfr1{10}$ program (10) has $465$ inequalities, of which $Ms(N_t+1)=3\cdot7\cdot13=273$ are safety rows and $p_UQ=6\cdot32=192$ are input-facet rows; the safety-row matrix has full column rank $r=mQ=64$, so Theorem~3.2(v) allows a positive certificate with up to $65$ safety labels. The dual witness of Proposition~5.3 uses three: the horizon labels $(j,\text{critical facet},T)$, $j=1,2,3$, with weights $\lambda=(\tfr38,\tfr5{16},\tfr5{16})$, plus input-facet multipliers $\mu_{j,k}=\lambda_ja_k\eta_j$ on the post-observation blocks (these are not safety labels and do not count toward the support). The deletion-based procedure of Section~8.1 run on this certificate with target margin $\gamma_0=\tfr3{100}$ attempts each of the three supported labels in turn and keeps none of the deletions: removing the label of branch $j$ leaves the labels of the complementary pair, and the pair's finite program is feasible (its optimum is at most the pair's exact slack, $-\tfr{31}{1250}$ or $-\tfr{81}{1250}$, below zero), so no certificate---let alone one with the target margin---survives. The support $\{1,2,3\}$ is therefore irredundant, and by Proposition~4.3 it is minimal, since the obstruction depth of the triple at $\tau=\tfr15$ is $3$. The gap between $3$ and the admissible $65$ is the statement of Corollary~5.5 in computational form: the rank bound is a worst case over all row systems of the given dimensions, and the depth is the instance's own conflict structure.

The same conclusion holds on every aligned mesh with $h<\tfr15$---for instance on the $2^7$-fold refinement of $\{0,\tau,T\}$ with $8469$ rows and up to $mQ+1=1025$ admissible labels---because the witness pair is mesh-generic and the pairs' programs remain feasible: the certificate's support does not grow with the mesh, only the program does. This is the practical content of the size result: on a nonviable instance one may expect the \emph{certificate} to stay small as the \emph{program} is refined, and the residual-bounded test (18) checks it in time linear in the support.
